\section{Discussion}

The frozen results show that feature fusion can help the HAM10000 benchmark when complementary backbone representations are kept available to the classifier. ResNet50 + EfficientNetB0 concatenation was the strongest frozen model, suggesting that the ResNet50 and EfficientNetB0 feature spaces captured useful non-identical information. Adding MobileNetV2 as a third concatenated backbone improved over the single-backbone baselines, but did not beat the best pairwise fusion.

Concatenation outperformed projected weighted fusion in this sprint. A likely explanation is that concatenation preserves the full feature vectors, while weighted fusion compresses each backbone into a shared 512-dimensional space before summation. The weighted models are still useful for interpretation: their learned weights did not collapse to a single backbone, and the three-backbone run assigned broadly distributed weights across all three CNNs.

Per-class F1 confirms why macro-F1 is more informative than accuracy. The best fusion model improved \texttt{df} F1 from 0.250 in the ResNet50 baseline to 0.480, and also improved \texttt{bcc}, \texttt{bkl}, \texttt{akiec}, and \texttt{mel}. These gains should be interpreted carefully because the minority-class supports are small, especially \texttt{df} and \texttt{vasc}. The work remains benchmark dermoscopic image classification and does not support clinical diagnosis claims.

The fine-tuned matrix shows that adapting the final CNN stages can improve the benchmark result
without changing the cached-feature MLP protocol. The best model shifted from frozen ResNet50 +
EfficientNetB0 concatenation to fine-tuned three-backbone concatenation, which suggests that
fine-tuning made the combined representation more useful to the downstream classifier. Weighted
fusion also benefited from fine-tuning, but concatenation remained the strongest final operator.

The main limitation of both the frozen and fine-tuned matrices is that they use one seed. The
fine-tuned gains are large enough to report, but they should not be treated as definitive
architecture rankings without additional seeds. This is especially important for minority classes
such as \texttt{df} and \texttt{vasc}, where small absolute changes in correct predictions can move
F1 substantially.

The representation complementarity analysis supports, but does not fully explain, the fusion results. The best concat model used the most complementary pair, ResNet50 + EfficientNetB0. However, the weighted model using the same pair did not achieve the same gain, which suggests that complementarity is useful only when the fusion operator preserves enough discriminative information. This strengthens the interpretation that concatenation was effective because it retained full backbone features, while projected weighted fusion traded some predictive capacity for compactness and interpretability.
